\chapter{Experiments}
\label{chap:ch4}

\par This chapter presents the experimental evaluation of the Digital Twin irrigation
    application described in Chapter~\ref{chap:ch3}. The purpose of the evaluation is to
    compare different irrigation strategies under consistent input conditions and to
    analyze their behavior in terms of water usage, scheduled irrigation events, soil
    moisture evolution, data efficiency, prediction quality, and actuation feasibility.

\par The selected strategies reflect the main directions identified in the literature
    review: rule-based threshold irrigation control, reduced sensor data acquisition,
    neuro-fuzzy irrigation estimation, and fuzzy Digital Twin decision making.
    Evapotranspiration and estimated evaporation are part of the Digital Twin soil
    moisture simulation, while the Fuzzy Digital Twin controller uses soil moisture,
    temperature, and rainfall as its fuzzy inputs for the resulting prescription.

\section{Dataset}
\label{sec:ch4sec1}

\par The dataset is composed of weather data, simulated or stored sensor readings, plant
    and pot configuration, sensor placement information, and irrigation records. The
    application was designed to support both historical analysis and forecast simulation.
    For each selected date range, data sources are loaded from PostgreSQL, generated from
    simulated sensor behavior when needed, and combined with estimated weather or Digital
    Twin state values into an experiment snapshot that provides a common input context for
    all irrigation strategies.

\par Weather data is stored at hourly granularity to provide the environmental context.
    The stored records include variables such as air temperature, relative humidity,
    precipitation, wind information, cloud cover, pressure fields, soil-related weather
    fields, shortwave radiation, and evapotranspiration when these values are available
    from the weather source. For future intervals, forecast weather is used when
    available; beyond the forecast horizon, the Digital Twin estimates missing weather and
    state values to maintain a continuous simulation interval. Weather ingestion is
    handled by the backend weather service, which imports Open-Meteo archive and forecast
    data and stores refresh audit records.

\par Sensor data represents soil moisture and local microclimate measurements for pot
    locations. Readings may be simulated to reproduce a physical sensing layer. The
    simulation follows state evolution rather than independent random generation: soil
    moisture evolves from the previous state according to evapotranspiration or estimated
    evaporation, precipitation, pot exposure to sun and wind, plant type, container
    properties, sensor calibration when readings exist, seasonal rules, and virtual
    irrigation effects. In the comparative ANFIS-GA-inspired and Fuzzy Digital Twin
    experiments, sensor readings are used to anchor the initial state, after which each
    strategy evolves soil moisture from weather effects and its own irrigation decisions.
    This avoids repeatedly pulling the experimental trajectories back toward the same
    stored sensor state during the simulated period. Sensor history is stored at tiered
    resolution: recent raw readings are kept at 15-minute resolution for the last 24
    hours, hourly readings are retained for the previous 7-day period, and older readings
    are stored for up to one year at configured daily slots to reduce storage
    requirements.

\par The plant and pot dataset defines plant type, pot size, balcony zone, rain exposure,
    sun exposure, wind exposure, drip flow, and moisture requirements. This allows
    controllers to apply plant-specific thresholds and valve-zone behavior instead of
    relying on a single global moisture rule. Since only selected pots are directly
    associated with sensors, the Digital Twin estimates the state of the remaining pots
    using weather data, pot properties, representative sensor readings, and irrigation
    effects from previous simulated or recorded watering events. In the simulated
    inventory, one representative sensor is assigned to each valve zone. The selected pot
    is treated as a zone sentinel and is chosen from the pots with the highest expected
    dry-down and water-need risk, such as small, exposed, heat-sensitive, or low-retention
    containers. This reflects the intended real-world placement: if the most vulnerable
    pot in a valve zone remains adequately moist, the other pots in the same valve zone
    are expected to be less likely to require earlier irrigation. If another pot
    repeatedly becomes the earlier stress indicator, the physical sensor would have to be
    moved in a later deployment; automatic sensor relocation is outside the current
    implementation scope. The controllers therefore use the selected sensor locations as
    the decision inputs. The remaining pots are still simulated so that the effect of
    valve-zone watering can be evaluated across the full inventory, but they do not act
    as independent irrigation triggers in the current control experiments.

\par The database also stores irrigation prescriptions, planned irrigation events, and
    actuation records. These records connect experiment outputs with the actuation layer
    and support the simulated closed-loop workflow. In the current implementation,
    detailed irrigation decisions and alerts are returned as experiment output for
    analysis, while the runtime actuation pipeline is configured to materialize and
    consume baseline prescriptions as simulated actuator work items.

\par When stochastic elements are involved, deterministic seeds are used to reproduce
    generated data, such as the pot inventory or synthetic environmental sequences. This
    ensures that differences in results are primarily caused by the irrigation strategy
    rather than by changing generated scenario data.

\section{Design of Experiments}
\label{sec:ch4sec2}

\par The experimental design follows a comparative approach, where all strategies are
    evaluated over the same date range, pot inventory, plant configuration, weather
    context, sensor placement assumptions, and experiment snapshot. The snapshot combines
    weather records, selected weather rows, day profiles, sensor context, initial pot
    states, plant moisture requirements, and selected sensor locations. For future
    intervals, it may also include estimated weather and simulated Digital Twin state
    evolution. When exact historical sensor readings are unavailable, the system uses
    retained sensor context and maps dates to comparable sensor records when possible. The
    representative sensor placement is therefore an input assumption shared by the
    compared strategies, not a dynamic controller decision during the experiment.
    Irrigation decisions are made from the selected sensor locations only. When a sensor
    location triggers a valve zone, the simulated valve delivery is applied to all active
    pots in that zone according to their drip flow rates. The plotted moisture values are
    aggregate outcome metrics for the full simulated inventory, not the direct decision
    signal used by the controller. The baseline controller uses sensor calibration
    markers during the simulation, because it represents the active Digital Twin
    reference strategy. For active-season analyses, baseline state warm-up is limited to
    the beginning of the current irrigation season rather than the dormant winter period;
    dormant-season intervals start from the selected range itself. The reduced sampling
    experiment reuses the baseline-warmed pot states at the selected start date, then
    applies sparse sensor refreshes from that shared state. In contrast, the
    ANFIS-GA-inspired and Fuzzy Digital Twin simulations use this sensor context only as
    an initial state anchor and then evolve their moisture trajectory from weather
    effects and their own irrigation decisions. This distinction means that the baseline
    is a calibrated operational reference, reduced sampling is a data-frequency variant
    of that reference state, and ANFIS-GA-inspired and fuzzy results are counterfactual
    strategy simulations after the same initial anchoring step.

\par The design of experiments for the Digital Twin irrigation evaluation is represented
    in Figure~\ref{designOfExperiments}. The same experiment snapshot is used to compare
    the baseline controller, reduced sensor sampling, ANFIS-GA-inspired irrigation signal,
    and Fuzzy Digital Twin prescription strategy through shared metrics such as water
    usage, moisture stability, irrigation events, and valve runs.

\begin{figure}[htbp]
	\centering
        \includegraphics[scale=0.2]{./figures/designOfExperiments}
	\caption{Design of experiments for the Digital Twin irrigation evaluation}
	\label{designOfExperiments}
\end{figure}

\par One reference strategy and three experimental variants were implemented. The
    reference strategy is the baseline Digital Twin irrigation controller. Although it is
    used as the comparison point, it is not a simple static threshold rule. At predefined
    irrigation decision slots, it evaluates irrigation need through dynamic thresholds. If
    irrigation is required, it calculates water volume from the current moisture deficit.
    The valve runtime is requested by the selected sensor location, while delivery is
    later distributed to the active pots in the same valve zone. The compact calculation
    is shown in Listing~\ref{lst:water_need_calculation}.

\begin{lstlisting}[
    language=Python,
    caption={Sensor-based irrigation request calculation},
    captionpos=b,
    label={lst:water_need_calculation},
    breaklines=true,
    basicstyle=\scriptsize\ttfamily
]
request_sensor_id = int(decision["sensor_id"])
need_pct = max(0.0, target - start_moisture)
flow_rate_ml_min = max(pot["drip_flow_ml_min"] * flow_multiplier, 1.0)

full_dose_ml = min(
    need_pct * pot["volume_l"] * 10.0 / retention,
    flow_rate_ml_min * max_minutes,
)
planned_ml = full_dose_ml * dose_factor
duration_min = planned_ml / flow_rate_ml_min

return {
    "request_sensor_id": request_sensor_id,
    "planned_volume_ml": round(planned_ml, 2),
    "duration_min": round(duration_min, 3),
}
\end{lstlisting}

\par Simulated sensor generation is a separate process used to populate the sensor
    readings table. The baseline experiment can use these readings to anchor or calibrate
    the Digital Twin state, but the irrigation simulation itself still evolves soil
    moisture and computes irrigation volume during the experiment.

\par The second experiment evaluates reduced sensor sampling. Instead of assuming that
    direct sensor readings are available continuously, the experiment reduces the
    frequency at which the Digital Twin receives sensor updates. Between direct refreshes,
    soil moisture is estimated using the latest known state, weather data, pot properties,
    and simulated irrigation effects. To avoid a second independent pre-simulation, the
    sparse strategy starts from the same baseline-warmed pot states at the beginning of
    the selected interval. This experiment is relevant for IoT irrigation systems because
    reducing sensor acquisition frequency can reduce communication overhead, storage
    volume, and energy consumption at the sensing layer.

\par The third experiment is inspired by an Adaptive Neuro-Fuzzy Inference System
    optimized with a genetic algorithm. In the implemented version, a compact ANFIS-style
    model estimates irrigation need using soil moisture, temperature, and effective rain.
    The model is trained from stored sensor-weather examples, with additional weighting
    for dry, hot, and post-irrigation recovery situations. The output is interpreted as an
    irrigation probability. During execution, the probability is evaluated at valve-zone
    level and is combined with managed-valve eligibility, moisture safety thresholds,
    water-saving cadence logic, and ANFIS-specific dose scaling based on probability
    confidence and moisture deficit. Rainfall and temperature influence the ANFIS decision
    through the model inputs rather than through the same explicit rain-aware dose rules
    used by the baseline controller.

\par The fourth experiment is a rule-based Fuzzy Digital Twin irrigation prescription
    strategy. It uses soil moisture, temperature, and rain amount as primary fuzzy inputs,
    which are combined through fuzzy rules to produce the irrigation prescription. Plant
    and pot requirements influence the safety references and the conversion of the
    prescription into water volume. The fuzzy prescription is used as the main activation
    signal, while fixed rules are limited to physical and operational guards such as cold
    conditions, irrigation windows, and valve-zone execution. This experiment evaluates
    whether a fuzzy decision layer can adapt irrigation prescriptions to current
    environmental conditions and produce different water-use behavior from the baseline
    strategy.

\par In the comparative experiments, the sampling, ANFIS-GA-inspired, and fuzzy strategies
    generate simulated irrigation decisions and valve-zone events. However, the runtime
    actuation pipeline is configured to materialize and consume only the baseline
    prescription. Therefore, the experimental results are used for evaluation and
    comparison, while the baseline remains the active control strategy for the simulated
    actuator dispatch workflow.

\par The experiments are not intended to reproduce a complete real agricultural
    deployment. Instead, they provide a controlled environment for comparing irrigation
    strategies and validating the behavior of the developed application. This is
    consistent with the role of the Digital Twin as a virtual environment in which
    irrigation strategies and actuation plans can be evaluated before real-world
    deployment.

\section{Analysis Metrics}
\label{sec:ch4sec3}

\par Experiment results are evaluated using the metrics defined in Chapter~\ref{chap:ch3}.
    The main comparison criteria are total water usage, irrigation days, valve runs,
    scheduled irrigation events, soil moisture evolution, and moisture stability. These
    metrics show how each strategy affects water consumption and whether soil moisture
    remains close to a comfortable range.

\par Scheduled irrigation events indicate when the controller determines that watering is
    required within the allowed irrigation windows. In the implemented system, irrigation
    is not triggered automatically every day. Decisions depend on soil moisture,
    plant-specific thresholds, pot and valve-zone configuration, temperature, rainfall,
    forecast conditions, and the selected control strategy. An efficient strategy should
    therefore avoid unnecessary watering when soil moisture is sufficient or when rainfall
    can contribute to the moisture requirement.

\par For the reduced sampling experiment, the analysis focuses on the effect of lowering
    direct sensor acquisition frequency. The reduced sampling strategy is compared with
    the baseline using water usage, decision agreement, mismatch days, valve runs, missed
    and extra valve activations, moisture estimation error, estimation bias, maximum
    error, and refresh counts. This indicates whether fewer sensor updates can still
    support acceptable Digital Twin state estimation and irrigation behavior.

\par For the ANFIS-GA-inspired experiment, the analysis combines model fit indicators with
    irrigation behavior metrics. The model fit indicators include RMSE, probability fit
    percentage, category accuracy, decision accuracy, validation sample count, decision
    threshold behavior, and predicted irrigation probability range. These values are
    interpreted together with water usage, irrigation events, valve runs, moisture-safe
    savings, and soil moisture evolution, because a good probability model is only useful
    if it produces reasonable irrigation decisions and preserves acceptable moisture
    conditions.

\par For the Fuzzy Digital Twin experiment, the analysis considers fuzzy prescription
    values, applied water, irrigation days, valve-zone behavior, comfort-preserved
    savings, moisture-safe savings, and the resulting soil moisture trajectory. The fuzzy
    controller uses soil moisture, temperature, and rain amount as its fuzzy inputs.

\par Since the application also includes actuation logic, the results can be interpreted
    using valve-level indicators such as actuator runtime estimates, number of valve runs,
    planned volume, and valve-zone scheduling feasibility. The overview layer also
    computes an optimized valve schedule, maximum parallel flow, and whether the planned
    sequence fits within the next available irrigation window. These values support
    feasibility analysis, but they should be interpreted as scheduling heuristics rather
    than as a complete hydraulic validation model.

\section{Results}
\label{sec:ch4sec4}

\par This section presents the results obtained after running the implemented experiments.
    The results are analyzed using time-series visualizations showing soil moisture
    evolution, irrigation events, rainfall, temperature, water usage, and, where
    applicable, predicted values for future intervals.

\par Visualizations distinguish historical or simulated values from forecast values,
    clarifying which part of the chart is supported by stored data and which part relies
    on forecast weather and Digital Twin state estimation. The comfort threshold displayed
    as a chart reference line corresponds to the average target moisture of the active
    pots, rather than the minimum admissible moisture or the exact irrigation trigger used
    by the controller. For the daily views used in the figures, the plotted moisture
    represents the end-of-day value, so this threshold should be interpreted as a
    conservative comfort benchmark: values below the line do not necessarily indicate
    plant stress, but indicate that the substrate ended the day below the preferred refill
    target.

\subsection{Baseline Threshold Experiment}
\label{baseline_results}

\par In the baseline experiment, irrigation is scheduled by a deterministic threshold and
    weather-aware controller derived from each pot's moisture requirements. The trigger
    threshold and dose are adjusted using contextual weather information, including
    temperature, precipitation amount and probability, pot rain exposure, dry and windy
    conditions, dry streaks, and dormant season constraints.

\par Figure~\ref{BaselineExperiment} illustrates the baseline end-of-day soil moisture
    evolution together with water usage, rainfall, and maximum temperature. The plotted
    moisture therefore represents the final simulated moisture state at the end of each
    day, after the daily effects of irrigation, rainfall, and environmental moisture loss
    have been applied, rather than a measurement taken immediately after irrigation.

\begin{figure}[htbp]
	\centering
        \includegraphics[scale=0.62]{./figures/baseline}
	\caption{Baseline threshold experiment showing soil moisture evolution, water usage,
	    rainfall, and maximum temperature. The figure illustrates how rainfall contributes to
	    moisture recovery and how irrigation volume is reduced during rainy periods, while
	    sheltered pots may still require watering.}
	\label{BaselineExperiment}
\end{figure}

\par The figure shows that rainfall has a visible effect on the simulated soil moisture
    trajectory, especially during the forecast interval, where soil moisture increases
    after rainy periods. At the same time, irrigation volume is often reduced during rainy
    days, indicating that the controller accounts for precipitation contribution.
    Irrigation does not always drop to zero because not all pots are equally exposed to
    rain; sheltered pots may still require watering, although with a reduced dose.

\par Overall, the baseline provides a realistic reference behavior for the system. It
    reacts to low soil moisture, incorporates weather context, and respects the irrigation
    window logic. More advanced strategies are therefore evaluated against this reference
    to determine whether they can reduce water usage, maintain moisture stability,
    decrease unnecessary irrigation events, or improve decision quality.

\subsection{Reduced Sampling Experiment}

\par The reduced sampling experiment evaluates the impact of lowering the frequency of
    sensor data acquisition. It shows how irrigation behavior changes when the Digital
    Twin receives fewer direct sensor updates and relies more heavily on estimated soil
    moisture evolution between measurements.

\par A moderate reduction in sampling frequency can still preserve acceptable irrigation
    behavior if the Digital Twin model estimates moisture accurately. If the sampling
    interval becomes too large, the system may react later to changes in soil moisture,
    leading to delayed irrigation, lower accuracy, or increased deviation from the
    baseline.

\par Table~\ref{sparse_sensing_comparison} summarizes the comparison between the 48 h and
    72 h sampling configurations.

\begin{table}[htbp]
\centering
\caption{Comparison of sparse sensing experiment results}
\label{sparse_sensing_comparison}
\begin{tabular}{lcc}
\toprule
\textbf{Metric} & \textbf{48 h sampling} & \textbf{72 h sampling} \\
\midrule
Interpretation & Good trade-off & Needs review \\
Water saved & 15.37\% & 6.73\% \\
Water saved vs. baseline & 339.75 L & 148.71 L \\
Decision agreement & 100.0\% & 96.5\% \\
Mismatch days & 0 & 1 \\
Sparse / baseline valve activations & 111 / 102 & 113 / 102 \\
Valve activation delta & +9 & +11 \\
Valve shift & 8.8\% & 10.8\% \\
Missed / extra activations & 0 / 9 & 0 / 11 \\
\bottomrule
\end{tabular}
\end{table}

\par The 48 h sampling configuration produced the strongest result, preserving full
    decision agreement with the baseline. It reached 100.0\% agreement with no mismatch
    days, while reducing water use by 339.75 L, equivalent to 15.37\% savings. It produced
    111 valve activations compared with 102 in the baseline, resulting in an 8.8\% valve
    shift. Since no baseline activations were missed, the 48 h interval represents a good
    trade-off between reduced sensing frequency and irrigation-control quality.

\par The 72 h configuration also reduced water use, but with weaker alignment. It saved
    148.71 L, equivalent to 6.73\% savings, and achieved 96.5\% decision agreement with
    one mismatch day. Valve activations increased from 102 in the baseline to 113,
    resulting in a 10.8\% valve shift. Although no baseline irrigation activations were
    missed, the lower savings and higher shift indicate that the 72 h interval should be
    treated as acceptable but requiring review.

\par Figure~\ref{SamplingExperiment} compares the two sampling intervals in terms of
    moisture, irrigation volume, and weather-aware decision behavior. The 48 h case stays
    closer to the baseline, while the 72 h case introduces more deviation and therefore
    requires additional review despite still reducing water consumption.

\begin{figure}[htbp]
	\centering
		\includegraphics[scale=0.43]{./figures/sampling48h}
        \includegraphics[scale=0.43]{./figures/sampling72h}
	\caption{Comparison of sparse sensing for 48 h and 72 h sampling intervals}
	\label{SamplingExperiment}
\end{figure}

\par Figure~\ref{SamplingRainForecast} shows that the sparse sensing trajectory follows
    the baseline moisture evolution very closely after rainfall, but volumes differ due to
    different state estimates.

\begin{figure}[htbp]
	\centering
	\includegraphics[scale=0.62]{./figures/SamplingRainForecast.png}
	\caption{Rainfall effect during forecast simulation}
	\label{SamplingRainForecast}
\end{figure}

\par This indicates that the Digital Twin can estimate moisture recovery after rain even
    when direct sensor readings are sparse. However, the irrigation volume is not
    identical in the two strategies. Although the resulting end-of-day moisture values are
    similar, the baseline and sparse sensing strategies may apply different water
    quantities because they rely on different sensor corrections. As a result, the same
    weather context can lead to different irrigation volumes. Therefore, similar moisture
    trajectories do not necessarily imply identical irrigation actions.

\par It should also be noted that this example was generated after the weather forecast
    had been refreshed. Since rainfall forecasts can change significantly from one day to
    another, the later simulation is not directly identical to the earlier 48 h and 72 h
    comparison. This highlights an important limitation: future results depend on the
    forecast version available at execution time.

\par These results show that sparse sensing can reduce the need for frequent sensor
    acquisition while preserving irrigation decisions, especially at the 48 h interval.
    Increasing the interval to 72 h still maintains high agreement, but the Digital Twin
    relies more heavily on estimated soil moisture between measurements, which increases
    the risk of shifted irrigation timing. The additional rainfall example shows that
    sparse sensing can also reproduce the baseline moisture trajectory after
    precipitation, even when the applied irrigation volume differs. From a distributed
    systems perspective, reducing continuous sensing by estimating intermediate states can
    reduce data traffic, storage volume, and energy consumption at the sensing layer.

\subsection{ANFIS-GA-Inspired Experiment}
\label{subsec:anfis_results}

\par This experiment evaluates whether irrigation can be supported by a learned
    probability signal instead of deterministic moisture thresholds alone. The
    ANFIS-GA-inspired model estimates irrigation need from soil moisture, temperature,
    and effective rain. The selected range is anchored to the same initial Digital Twin
    state as the baseline, after which the ANFIS moisture trajectory evolves from its own
    decisions, weather effects, and simulated irrigation delivery.

\par The ANFIS output is not a direct actuator command. It is interpreted at valve-zone
    level and combined with moisture eligibility checks, safety thresholds, managed-valve
    constraints, cadence logic, and dose scaling. Therefore, three metrics are separated
    in the interpretation: baseline agreement describes runtime agreement with the
    threshold controller; validation agreement describes binary agreement with the ANFIS
    training target; and target probability fit describes the continuous probability fit
    to that target.

\begin{table}[htbp]
\centering
\scriptsize
\caption{ANFIS-GA-inspired experiment metrics}
\label{tab:anfis_results_metrics}
\begin{tabular}{|p{2.7cm}|p{2.0cm}|p{2.0cm}|p{1.7cm}|p{1.8cm}|p{2.1cm}|}
\hline
\textbf{Interval} & \textbf{Baseline agreement} & \textbf{Validation / fit} &
\textbf{Valve shift} & \textbf{Water result} & \textbf{Execution core / API} \\
\hline
29.05.2026--26.06.2026 & 93.10\% (2 days) & 93.01\% / 91.54\% &
$-4/+35$ & 126.52 L extra & 1.00 s / 1.65 s \\
\hline
29.05.2023--26.06.2023 & 72.41\% (8 days) & 91.68\% / 90.16\% &
$-4/+39$ & 273.46 L extra & 1.14 s / 1.70 s \\
\hline
29.01.2026--26.02.2026 & 72.41\% (8 days) & 93.01\% / 91.54\% &
$0/+19$ & 102.20 L extra & 1.20 s / 3.01 s \\
\hline
\end{tabular}
\end{table}

\par Figure~\ref{fig:anfis_results2} presents the current-season run. The baseline
    agreement was 93.10\%, with two decision days different from the baseline. The
    validation indicators were also high: 93.01\% validation agreement and 91.54\%
    target probability fit over 8170 evaluation samples. The ANFIS run used more valve
    activations than the baseline, producing 126.52 L of additional water use in this
    execution, while preserving comfort on 27 of 29 days.

\begin{figure}[htbp]
	\centering
    \includegraphics[scale=0.5]{figures/anfis2.PNG}
	\caption{ANFIS-GA-inspired irrigation strategy evaluation showing baseline moisture,
	    ANFIS estimated moisture, ANFIS signal, irrigation volume, and weather context}
	\label{fig:anfis_results2}
\end{figure}

\par The historical seasonal run is shown in Figure~\ref{fig:anfis_results2023}. Since
    exact 2023 sensor rows are not available, the interval should be interpreted as a
    comparable stored-weather scenario rather than a replay of measured 2023 plant
    state. The 2023-specific ANFIS model was selected automatically from the lower date
    bound. Its validation agreement remained 91.68\%, with 90.16\% target probability
    fit over 3460 evaluation samples. Runtime baseline agreement was lower, at 72.41\%,
    because the learned signal activated more often than the conservative historical
    baseline. The additional 273.46 L should therefore be read as comfort-preserving
    model behavior under cooler stored-weather conditions, not as an isolated failure of
    the probability model.

\begin{figure}[htbp]
	\centering
    \includegraphics[scale=0.5]{figures/anfis2023.PNG}
	\caption{ANFIS-GA-inspired irrigation strategy evaluation for a historical seasonal
	    interval, showing that lower temperature and different rainfall timing can reduce
	    irrigation demand, even when rainfall is not directly comparable with the 2026
	    interval.}
	\label{fig:anfis_results2023}
\end{figure}

\par The dormant-season result is shown in Figure~\ref{fig:anfisDormantSeason}. During
    December to March, the baseline applies a strict winter rule and irrigates only under
    narrow temperature, rain, and moisture conditions. ANFIS does not reuse that fixed
    winter skip policy; it evaluates the learned valve-zone signal together with moisture
    deficit and safety constraints. For 29.01.2026--26.02.2026, the baseline used no
    water, while ANFIS scheduled 19 extra valve activations on eight days and applied
    102.20 L. A percentage water-saving value is not meaningful against a zero-water
    baseline. The 72.41\% baseline agreement comes from the fact that most dormant days still
    matched as no-irrigation days.

\begin{figure}[htbp]
	\centering
    \includegraphics[scale=0.52]{figures/anfisDormantSeason.PNG}
	\caption{ANFIS-GA-inspired behavior during dormant season without the same fixed winter
	    skip rule as the baseline}
	\label{fig:anfisDormantSeason}
\end{figure}

\par The ANFIS endpoint loads a persisted model; training is performed separately and is
    not part of the request time. The reported backend execution time ranged from 1.00 s
    to 1.20 s for the tested intervals, while measured end-to-end API time ranged from
    1.65 s to 3.01 s. The API time also includes snapshot preparation, weather access,
    persistence, JSON serialization, and transfer of the full experiment payload. The
    current-season request is expected to be heavier because it can include forecast and
    estimated-weather handling.

\subsection{Fuzzy Digital Twin Prescription Experiment}
\label{subsec:fuzzy_results}

\par The Fuzzy Digital Twin experiment compares an adaptive fuzzy irrigation prescription
    strategy with the default threshold controller. The fuzzy controller estimates an
    irrigation prescription using three primary fuzzy inputs: soil moisture, temperature,
    and rainfall. These inputs are combined through fuzzy rules to produce a prescribed
    irrigation depth, expressed in millimetres. The resulting depth is converted into an
    irrigation volume according to the pot surface area. As with the ANFIS-GA-inspired
    experiment, the fuzzy simulation uses sensor readings as an initial state anchor and
    then evolves its own moisture trajectory from fuzzy prescriptions, weather effects,
    and simulated irrigation delivery.

\begin{table}[htbp]
\centering
\scriptsize
\caption{Fuzzy Digital Twin prescription metrics}
\label{tab:fuzzy_results_metrics}
\begin{tabular}{|p{2.7cm}|p{2.0cm}|p{1.7cm}|p{1.8cm}|p{1.8cm}|p{1.8cm}|p{2.1cm}|}
\hline
\textbf{Interval} & \textbf{Baseline agreement} & \textbf{Valve shift} &
\textbf{Prescription} & \textbf{Water result} & \textbf{Comfort} &
\textbf{Execution core / API} \\
\hline
29.05.2026--26.06.2026 & 93.10\% (2 days) & $0/+85$ & 47.72 L/day avg &
315.41 L saved & 12/29 days & 1.29 s / 2.00 s \\
\hline
29.05.2023--26.06.2023 & 79.31\% (6 days) & $0/+102$ & 33.83 L/day avg &
206.02 L extra & 29/29 days & 1.23 s / 1.89 s \\
\hline
29.01.2026--26.02.2026 & 68.97\% (9 days) & $0/+43$ & 6.56 L/day avg &
190.16 L extra & 29/29 days & 1.08 s / 1.39 s \\
\hline
\end{tabular}
\end{table}

\par Figure~\ref{fig:fuzzyCurrent} shows the current-season fuzzy prescription run. The
    strategy agreed with the baseline on 93.10\% of days and saved 315.41 L, but it did so
    by using many small valve-zone prescriptions: 151 fuzzy valve activations compared
    with 66 baseline activations. Comfort was preserved on 12 of 29 days, so the saved
    water must be interpreted together with lower moisture comfort rather than as an
    isolated optimization result. The average prescription was 47.72 L/day.

\begin{figure}[htbp]
	\centering
    \includegraphics[scale=0.52]{figures/fuzzyCurrent.PNG}
	\caption{Fuzzy Digital Twin experiment for the current period}
	\label{fig:fuzzyCurrent}
\end{figure}

\par The historical seasonal interval is shown in Figure~\ref{fig:fuzzyHistorical}. As in
    the ANFIS historical run, the absence of exact 2023 sensor rows means that the result
    should be read as a stored-weather scenario. The fuzzy controller used 206.02 L more
    water than the baseline because it activated more frequently: 141 valve activations
    compared with 39 baseline activations. The individual prescriptions remained moderate,
    but the number of activations outweighed the smaller dose per activation. The result
    preserved comfort on all 29 days, which explains why extra water can appear even when
    the prescription logic is responding to moisture and weather rather than overwatering
    indiscriminately.

\begin{figure}[htbp]
	\centering
    \includegraphics[scale=0.52]{figures/fuzzyHistorical.PNG}
	\caption{Fuzzy Digital Twin prescription for a historical seasonal interval}
	\label{fig:fuzzyHistorical}
\end{figure}

\par The dormant-season result is presented in Figure~\ref{fig:fuzzyDormant}. The baseline
    did not irrigate because the winter rule prevented watering under the selected
    conditions. The fuzzy controller does not apply that fixed winter skip rule; it still
    evaluates low moisture, temperature, and rainfall through its prescription rules. The
    average prescription was only 6.56 L/day, but it produced 43 extra valve activations over
    nine days and 190.16 L of water. As with ANFIS, a percentage increase is not meaningful
    because the baseline volume was zero. The 68.97\% baseline agreement reflects that
    most dormant days still matched the baseline no-irrigation decision.

\begin{figure}[htbp]
	\centering
    \includegraphics[scale=0.52]{figures/fuzzyDormant.PNG}
	\caption{Fuzzy Digital Twin behavior during dormant season without the same fixed winter
	    skip rule as the baseline}
	\label{fig:fuzzyDormant}
\end{figure}

\par The fuzzy endpoint is rule-based and does not require model loading or training. The
    reported backend execution time ranged from 1.08 s to 1.29 s, while measured
    end-to-end API time ranged from 1.39 s to 2.00 s. The difference between core and API
    time comes from the same surrounding work as the other experiment endpoints:
    snapshot preparation, weather access, persistence, serialization, and response
    transfer.

\section{Comparison of Irrigation Strategies}
\label{sec:ch4sec5}

\par The implemented experiments differ in complexity, required data, and decision-making
    capability. The baseline is a deterministic threshold and weather-aware Digital Twin
    strategy. Its main advantage is interpretability, because the irrigation action can be
    traced to moisture thresholds, weather modifiers, rain rules, and valve-zone
    constraints. However, it remains rule-based and depends on manually designed
    thresholds and dose factors.

\par The reduced sampling experiment does not introduce a new irrigation rule, but instead
    evaluates the effect of limited sensor availability. Its value lies in measuring how
    much data acquisition can be reduced while still maintaining acceptable control
    behavior. This experiment is relevant for scalable IoT and Digital Twin systems, where
    sensors may be battery-powered or communication may be constrained.

\par The ANFIS-GA-inspired experiment introduces a learned neuro-fuzzy decision-support
    layer. It can model nonlinear relationships between moisture state, weather context,
    and irrigation need. Compared with the threshold baseline, it can potentially adapt
    better to combined environmental conditions, but it requires training data, persisted
    model calibration, and runtime execution constraints that translate the learned signal
    into valve-zone actions.

\par The Fuzzy Digital Twin experiment introduces an explicit fuzzy prescription layer. It
    translates soil moisture, temperature, and rainfall into an interpretable irrigation
    depth. Its main value is transparent prescription control; water use is one outcome
    and must be interpreted together with moisture comfort, valve-zone actuation, and
    seasonal conditions.

\par Comparing the two intelligent strategies, ANFIS provides learned probability
    estimates with validation agreement and target probability fit, while the fuzzy
    controller provides explicit rule-based prescription values without a training step.
    In the current-season interval, both reached 93.10\% baseline agreement. In the
    historical interval, ANFIS reached 72.41\% agreement and fuzzy control reached
    79.31\% agreement because both activated more often than the conservative baseline.
    During the dormant interval, ANFIS reached 72.41\% agreement and fuzzy control
    reached 68.97\%, with both deviations caused by
    model-only irrigation days while the baseline winter rule skipped irrigation. In the
    measured runs, backend core execution stayed below 2 s for both strategies; end-to-end
    API time was also affected by snapshot preparation, weather access, persistence, and
    payload serialization, so these timings should be read as implementation measurements
    rather than hardware-independent benchmarks.

\par Table~\ref{tab:experiment_comparison} summarizes the main characteristics of the
    implemented strategies.

\begin{table}[h]
\centering
\caption{Comparison of implemented irrigation strategies}
\label{tab:experiment_comparison}
\begin{tabular}{|p{3.2cm}|p{3.2cm}|p{3.2cm}|p{3.2cm}|}
\hline
\textbf{Strategy} & \textbf{Main Inputs} & \textbf{Main Advantage} & \textbf{Main
Limitation} \\
\hline
Threshold baseline & Soil moisture, plant thresholds, weather and rain context, pot and
valve configuration & Deterministic and interpretable reference & Rule tuning is required
and adaptability is limited by predefined logic \\
\hline
Reduced sampling & Sparse sensor readings, Digital Twin state estimation, weather context,
pot associations & Reduces sensing and communication requirements & May shift irrigation
timing when sampling is too sparse \\
\hline
ANFIS-GA-inspired & Soil moisture, temperature, effective rain & Captures nonlinear
relations and provides a learned irrigation probability signal & Requires training,
persisted calibration, and additional execution policy to translate probabilities into
valve actions \\
\hline
Fuzzy Digital Twin & Soil moisture, temperature, rainfall, fuzzy rules & Explainable
prescription logic and direct dose control & May create frequent small activations; the
comfort and water outcome depends on rule tuning \\
\hline
\end{tabular}
\end{table}

\par Overall, the comparison shows that each strategy has a different role. The baseline
    provides reference behavior, the sampling experiment evaluates data efficiency, the
    ANFIS-GA-inspired method demonstrates learned irrigation-need estimation, and the
    Fuzzy Digital Twin experiment evaluates adaptive prescription using current and
    forecasted context. When prescriptions are dispatched, the actuation layer can
    translate irrigation requirements into valve runtimes and scheduled irrigation work
    items, but the current simulated actuator pipeline is configured for baseline
    prescriptions.

\section{Threats to Validity}
\label{sec:ch4sec6}

\par Several threats to validity must be considered when interpreting the results. The
    first limitation concerns the use of simulated sensor data. Although the simulation
    models soil moisture evolution using weather, evapotranspiration or estimated
    evaporation, rainfall, pot exposure, plant type, and container-related factors, it
    cannot fully reproduce all physical processes that occur in real soil and plant
    systems. Factors such as root depth, drainage behavior, detailed soil composition,
    plant growth stage, and local microclimate variation are simplified in the current
    implementation.

\par A second threat is related to the use of forecast weather data. Forecasts are
    inherently uncertain, and future weather conditions may differ from the values
    predicted at execution time. Since the Digital Twin uses forecast data to estimate
    future soil moisture, changes in predicted rainfall, temperature, or humidity may
    affect irrigation recommendations and simulated water usage. Moreover, forecasts can
    change significantly from one day to another, so running the same experiment today may
    produce different results from running it yesterday, even if the selected date range
    is identical. This limitation applies only to records marked as forecast data from the
    weather API; stored historical weather records remain stable unless they are
    explicitly refreshed or replaced.

\par A third limitation concerns the representativeness of evapotranspiration and
    moisture-loss estimation. Although the application uses evapotranspiration values when
    available and estimates evaporation from weather variables when needed, these values
    are combined with pot-level modifiers rather than direct microclimate measurements for
    each pot. Weather API evapotranspiration may differ from the balcony microclimate,
    because shading, wind exposure, reflected heat, rain shelter, and container material
    affect drying. Therefore, evapotranspiration is treated as an approximate component of
    Digital Twin state evolution, not as measured ground truth and not as a direct fuzzy
    or ANFIS input.

\par A fourth limitation concerns actuation realism. The implemented actuator workflow
    stores and consumes simulated actuation records, but no real MQTT, LoRa, HTTP, or
    serial hardware bridge is implemented for physical valves or sensor nodes. In
    addition, the current scheduling logic considers simplified valve flow constraints,
    but does not model pump pressure, pipe losses, emitter clogging, or detailed hydraulic
    interactions between simultaneous valve runs.

\par Another limitation concerns deployment scale mainly through data volume and
    operational validation. The system already estimates the state of 200 pots using a
    limited number of representative sensors, so larger scenarios do not necessarily
    require one sensor per plant. However, larger deployments would increase the amount of
    weather, sensor, plant, zone, and actuation data that must be ingested, processed,
    stored, and visualized. Larger deployments may also introduce more heterogeneous soil,
    exposure, hydraulic, and irrigation layout conditions. Therefore, the results validate
    the workflow in a controlled potted-plant scenario, while larger deployments would
    require further testing for data throughput, calibration, and robustness.

\par Another threat concerns the implementation of literature-inspired methods. The
    ANFIS-GA-inspired and Fuzzy Digital Twin experiments are adapted to the scope of the
    application. The ANFIS-GA-inspired controller learns an irrigation probability signal,
    but the runtime action is still shaped by Digital Twin execution constraints such as
    managed valve zones, safety thresholds, cadence logic, and dose scaling. The Fuzzy
    Digital Twin controller uses explicit fuzzy prescription rules rather than a learned
    model. They should therefore be interpreted as practical implementations inspired by
    the literature, rather than exact reproductions of the original systems.

\par The evaluation is also affected by parameter choices such as moisture thresholds, pot
    sizes, flow rates, sampling intervals, sensor placement, valve grouping, and fuzzy
    rules. Different parameter configurations may lead to different irrigation outcomes.
    For this reason, the application supports configurable experiments, but the
    conclusions should be interpreted in relation to the selected configuration.

\par Finally, the experiments evaluate software behavior in a controlled environment. A
    real deployment would require physical sensors, actuators, calibration, network
    communication, device failures, water pressure variation, and long-term plant health
    monitoring. These aspects are outside the current implementation scope but represent
    important directions for future work.

\section{Chapter Summary}
\label{sec:ch4summary}

\par This chapter presented the experimental evaluation of the Digital Twin irrigation
    application. The dataset combines weather data, simulated or stored sensor readings,
    plant-specific parameters, pot configurations, sensor placement data, and
    irrigation-related records. The experiments were designed to compare several
    irrigation strategies under consistent conditions by using a shared experiment
    snapshot.

\par The implemented experiments include baseline threshold and weather-aware irrigation,
    reduced sensor sampling, ANFIS-GA-inspired neuro-fuzzy irrigation estimation, and a
    Fuzzy Digital Twin prescription strategy. The evaluation focuses on water usage,
    scheduled irrigation events, soil moisture behavior, sensor data reduction,
    prediction-related behavior, valve runtime, and actuation scheduling feasibility.

\par The results show that the application can be used as an experimental environment for
    analyzing smart irrigation strategies. While the system does not replace real-world
    validation, it provides a useful virtual environment for testing Digital Twin
    irrigation logic, comparing decision strategies, and studying the trade-off between
    water efficiency, sensing requirements, moisture comfort, forecast sensitivity, and
    simulated actuation feasibility.
